\section{结合律公式}

\begin{proposition}{集合族乘积的``结合同构''}{}
	设$\left(X_{i}\right)_{i\in I}$是集合族,$I\neq\emptyset$,$\left(J_{\lambda}\right)_{\lambda\in\Lambda}$是$I$的划分,则
	\begin{align*}
		\phi:\prod\limits_{i\in I}X_{i}\to\prod\limits_{\lambda\in\Lambda}\left(\prod\limits_{i\in J_{\lambda}}X_{i}\right),f\mapsto\left(f|_{J_{\lambda}}\right)_{\lambda\in\Lambda}
	\end{align*}
	是双射.我们称$\phi$和$\phi^{-1}$为典范双射.
\end{proposition}

\begin{remark}{二元拆分的情形}{}
	如果$J=J_{\alpha}\sqcup J_{\beta}$,则存在典范双射$\prod\limits_{i\in I}X_{i}\to\left(\prod\limits_{i\in J_{\alpha}}X_{i}\right)\times\left(\prod\limits_{i\in J_{\beta}}X_{i}\right)$.特别地,若对任一$i\in J_{\beta}$,集合$X_{i}$是单点集,则部分投影$\pr_{J_{\alpha}}$是双射.
\end{remark}

\begin{remark}{三个集合的乘积和三者所构成集合族的乘积的关系}{}
	设$A,B,C$是集合,$I$由三个两两不相等的对象$\alpha,\beta,\gamma$构成,$X_{\alpha}=A,X_{\beta}=B,X_{\gamma}=C$,那么存在双射$\prod\limits_{i\in\left\{\alpha,\beta,\gamma\right\}}X_{i}\to A\times B\times C$.进一步的,下列图表交换,其中的每个箭头都是典范双射:
	\begin{center}
		\begin{tikzcd}
			& {A\times B\times C} & {A\times C\times B} & \\
			{B\times A\times C} &&& {C\times A\times B} \\
			& {B\times C\times A} & {C\times B\times A}
			\arrow["\simeq", no head, from=1-2, to=1-3]
			\arrow["\simeq", no head, from=1-3, to=2-4]
			\arrow["\simeq", no head, from=2-1, to=1-2]
			\arrow["\simeq"', no head, from=2-1, to=3-2]
			\arrow["\simeq", no head, from=2-4, to=3-3]
			\arrow["\simeq"', no head, from=3-2, to=3-3]
			\end{tikzcd}
	\end{center}
\end{remark}